% קובץ זה נוצר אוטומטית ע"י make_bank.py -- אל תערכו אותו ידנית.
% חוברת הרמזים והפתרונות לסטודנטים: רמזים (בלי פתרונות), ואחריהם התרגילים עם הפתרונות.
% מספרי התרגילים זהים למספרים בחוברת התרגילים (booklet.tex).
\newif\ifsol
\soltrue
\documentclass[12pt, a4paper]{article}
\usepackage[most]{tcolorbox}   % חייב להיטען לפני ה-preamble (bidi)
%\input{systeme}
\usepackage[a4paper, top=4cm,left=2.5cm, right=2.3cm, bottom=2.7cm]{geometry}
\usepackage{amsmath, amssymb, array, multicol}
\usepackage{comment, graphicx}  
\usepackage{pdfpages}
\usepackage{tikz}
\usepackage{fontspec}

\usepackage{pdfpages}
%\usepackage{xunicode}  % Hebrew character encoding
\usepackage{polyglossia}
\setmainlanguage{hebrew}
\usepackage{fontspec}
\setmainfont{David CLM}
\usepackage{bidi}
\usepackage{enumitem}

\usepackage{braket}
%\usepackage{xparse}
%\DeclareGraphicsExtensions{.png,.jpg}
%\usepackage{pgffor}  % for loops in latex
%\usepackage{ifthen}  % conditional latex compilation
%\usepackage{todonotes}  % notes in bright orange
%\usepackage[colorlinks=true]{hyperref}  % hyperlinks
%\AtBeginDocument{% coloring of links
%	\hypersetup{pdfborder={0 0 0},
%		urlcolor=blue}%
%}
\usepackage[absolute]{textpos}   % Proper positioning of images, e.g. BGU logo
%\setdefaultlanguage{hebrew}
%\setotherlanguage{english}
%\setmainfont{David CLM}
%\newfontfamily\hebrewfont[Script=Hebrew]{Times New Roman}
\setmainfont{David CLM}

%\makeatletter
%\def\tagform@#1{\maketag@@@{(\ignorespaces#1\unskip)}}
%\makeatother


%\makeatletter
%\renewcommand{\SepMark}[1]{\def\@SepMark{#1}}
%\makeatother
%\SepMark{.}
%\robustify\hebrewnumeral
%\robustify\Hebrewnumeral
%\robustify\Hebrewnumeralfinal

%\def\urladdrname{{\itshape כתובת רשת}}
%\def\emailaddrname{{\itshape כתובת דואל}}

%\makeatletter
%\let\orig@MakeUppercase \MakeUppercase
%\makeatother
%\providecommand{\MakeUppercase}[1]{#1}

\fontsize{14}{15}  

\setdefaultlanguage{hebrew}
\setotherlanguage{english}
\thispagestyle{empty}
\selectlanguage{hebrew}

\newcommand{\cc}{\mathbb{C}}
\newcommand{\CC}{\mathbb{C}}
\newcommand{\rr}{\mathbb{R}}
\newcommand{\RR}{\mathbb{R}}
\newcommand{\nn}{\mathbb{N}}
\newcommand{\zz}{\mathbb{Z}}
\newcommand{\ZZ}{\mathbb{Z}}
\newcommand{\qq}{\mathbb{Q}}
\newcommand{\FF}{\mathbb{F}}
\newcommand{\QQ}{\mathbb{Q}}
\DeclareMathOperator{\vol}{vol}
\newcommand{\abs}[1]{\left|#1\right|}

\DeclareMathOperator{\Adj}{Adj}
\DeclareMathOperator{\deter}{det}
\DeclareMathOperator{\Img}{Im}
\DeclareMathOperator{\Ker}{Ker}
\DeclareMathOperator{\rank}{rank}
\DeclareMathOperator{\Span}{Span}
\DeclareMathOperator{\tr}{tr}


\newcommand{\cB}{\mathcal{B}}
\newcommand{\cC}{\mathcal{C}}
\newcommand{\cR}{\mathcal{R}}
\newcommand{\cD}{\mathcal{D}}
\newcommand{\calE}{\mathcal{E}}
\newcommand{\ntn}{n\times n}
\newcommand{\mtn}{m\times n}
\newcommand{\limn}{\lim_{n\to\infty}}

%%%% new for complex functions

\newcommand{\Arg}{\mathrm{Arg}}
\newcommand{\ol}{\overline}
\newcommand{\wt}{\widetilde}
\newcommand{\wh}{\widehat}
\newcommand{\re}{\mathrm{Re}}
\newcommand{\im}{\mathrm{Im}}
\newcommand{\goes}{\longrightarrow}
\newcommand{\cO}{\mathcal{O}}
\newcommand{\Log}{\mathrm{Log}}



\def\I{\mathbb I}
\def\Q{\mathbb Q}
\def\F{\mathbb F}
\def\H{\mathbb H}
\def\HH{\mathcal H}
\def\PP{\mathcal P}
\def\cB{\mathcal B}
\def\cC{\mathcal C}
\def\cD{\mathcal D}
\def\cE{\mathcal E}
\def\mfp{\mathfrak p}
\def\mfm{\mathfrak m}
\def\GL{\operatorname{GL}}
\def\SL{\operatorname{SL}}
\def\Aut{\operatorname{Aut}}
\def\Af{\operatorname{Aff}}
\def\Gal{\operatorname{Gal}}
\def\Alt{\operatorname{Alt}}
\def\adj{\operatorname{Adj}}
\def\Sym{\operatorname{Sym}}
\def\Adj{\operatorname{Adj}}
\def\ord{\operatorname{ord}}
\def\sgn{\operatorname{sgn}}
\def\id{\operatorname{id}}
\def\tr{\operatorname{tr}}
\def\chr{\operatorname{char}}
\def\Ann{\operatorname{Ann}}
\def\Inn{\operatorname{Inn}}
\def\Hom{\operatorname{Hom}}
\def\Tor{\operatorname{Tor}}
\def\Rad{\operatorname{Rad}}
\def\Syl{\operatorname{Syl}}
\def\Hall{\operatorname{Hall}}
\def\ch{\operatorname{ch}}
\def\sp{\operatorname{Span}}
\def\Cay{\operatorname{Cay}}
\def\rank{\operatorname{rank}}
\def\cis{\operatorname{cis}}
\def\Geo{\operatorname{Geo}}
\def\Bin{\operatorname{Bin}}
\def\Ber{\operatorname{Ber}}
\def\sq{\subseteq}
\def\sqb{{\sss \char4}}
\def\emptyset{\mbox{\rm \O}}
\def\es{\emptyset}
\def\angle{<\!\!\!) \,\,}
\def\const{\mbox{\rm const}\,}
\def\lcm{\mbox{\L{\rm lcm}}\,}
\def\ga{\alpha}
\def\gb{\beta}
\def\gm{\gamma}
\def\gG{\Gamma}
\def\gf{\varphi}
\def\eps{\varepsilon}
\def\gd{\delta}
\def\gD{\Delta}
\def\s{\sigma}
\def\d{\partial}
\def\la{\lambda}
\def\La{\Lambda}
\def\Om{\Omega}
\def\om{\omega}
\def\gS{\Sigma}
\def\lenq{\lvertneqq}
\def\genq{\gvertneqq}
\def\ra{\Longrightarrow}
\def\cl{\centerline}
\def\ds{\displaystyle}
\def\ts{\dtextstyle}
\def\scr{\scriptsize}
\def\f{\frac}
\def\thefootnote{*}
\def\bs{\baselineskip}
\def\nbs{\normalbaselineskip}
\def\f5{{$\hskip 5pt$}}
\def\hqq{{\hskip 24pt}}
\def\hq{{\hskip 12pt}}
\def\longto{\mathop{\longrightarrow}\limits}
\def\mod{\mathop{\rm mod}\nolimits}
\def\cs{\choose}
\def\mbox{\hbox}
\def\se{\subseteq}
\def\qskip{\vskip 12pt}
\def\nl{\hfill\break}
\def\vbreak{\vfill\break}
\def\ol{\overline}
\def\sd{\triangle}
\def\sm{\smallsetminus}
\def\nin{\not\in}
\def\gg{\Gamma}
\def\omod#1{({\rm mod}\,\,#1)}
\def\qu{\qskip\item}
\def\bE{\begin{enumerate}}
\def\eE{\end{enumerate}}
\def\emq{\hskip 1in}
\def\p{\varphi}
\def\a{\forall}
\def\e{\operatorname{e}}
\def\z{\mathit{z}}
\def\h{\wedge}
\def\v{\vee}
\def\lcc{\langle}
\def\rcc{\rangle}
\def\gl{\mathbf {GL}}
\def\ker{\operatorname{ker}}
\def\mi{\operatorname{Im}}
\def\er{\operatorname{Re}}
\def\im{\mathbf{im}}
\def\norm{\trianglelefteq}
\newcommand{\Aro}{\Rightarrow}
\renewcommand{\atop}[2]{\genfrac{}{}{0pt}{1}{#1}{#2}}
\newcommand{\too}{\longrightarrow}
\renewcommand{\iff}[1]{\Longleftrightarrow}
\newcommand{\then}[1]{\hspace{#1}\Rightarrow\hspace{#1}}
\newcommand{\bcs}[2]{\genfrac{\{}{\}}{0pt}{1}{#1}{#2}}
\newcommand{\px}[1]{\begin{pmatrix}#1\end{pmatrix}}
\newcommand{\sx}[1]{\begin{bmatrix}#1\end{bmatrix}}
\newcommand{\dotcup}{\ensuremath{\mathaccent\cdot\cup}}
\newcommand{\cnt}[1]{\begin{center}#1\end{center}}
\newcommand{\cntl}[1]{\begin{center}\underline{#1}\end{center}}
\newcommand{\lfrac}[2]{\mathlarger{\frac{#1}{#2}}}
\newcommand{\aro}{\rightarrow}
\newcommand{\normal}{\trianglelefteq}
\newcommand{\spc}{\hspace{2 mm}}
\newcommand{\spcc}{\hspace{5 mm}}
\newcommand{\Spc}{\hspace{10 mm}}
\newcommand{\ff}{\varphi}
\newcommand{\cd}{\cdots}
\newcommand{\ml}[1]{\mathlarger{#1}}
\newcommand{\suml}[1]{\mathlarger{\sum #1}}
\newcommand{\intl}[1]{\mathlarger{\int #1}}
\newcommand{\norma}[1]{\abs{\abs{#1}}}
\newcommand{\lr}[1]{\left(#1\right)}
\newcommand{\lsr}[1]{\left[#1\right]}
\newcommand{\rseq}[3]{\left(#1_#2\right)_{#2=#3}^\infty}
\newcommand{\gen}[1]{\left\langle #1\right\rangle}
\newcommand{\seq}[3]{\gen{#1_#2}_{#2=#3}^\infty}
\newcommand{\floor}[1]{\left\lfloor#1\right\rfloor}
\newcommand{\case}[2]{\left\{\begin{array}{#1} #2 \end{array}\right.}
\newcommand{\bit}{\begin{itemize}}
\newcommand{\eit}{\end{itemize}}
\newcommand{\bM}[1]{\begin{multicols}{#1}}
\newcommand{\eM}{\end{multicols}}
\newcommand{\bT}[1]{\begin{tasks}(#1)}
\newcommand{\eT}{\end{tasks}}


\newcommand{\headers}{
\thispagestyle{empty}
\begin{textblock*}{9cm}[0.5,0.5](90mm,3cm)
	\includegraphics[width=2in]{preamble/logo.jpg}
\end{textblock*}	

\begin{center}
	\large	
היחידה להוראת קורסי יסוד - 
	\semester \\
	{\large \bf \course\ - 214.1.9921}\\
	\medskip
	{\Large  \textbf{\doctype}  }\\
\end{center}
\vspace{0.3cm}
}
%%%%%%%%%%%%%%%%%%%%%%%%%%%%%%%%%%%%%%%%%%%%%%%%%%%%%%%%%%%%%%%%%%%%%%
% הגדרות משותפות למאגר ולדפי התרגילים שנבנים ממנו.
% נטען אחרי preamble/MotkePreamble, ודורש ש-\ifsol יוגדר לפני \documentclass:
%     \newif\ifsol
%     \soltrue      % או \solfalse
%%%%%%%%%%%%%%%%%%%%%%%%%%%%%%%%%%%%%%%%%%%%%%%%%%%%%%%%%%%%%%%%%%%%%%
\tcbuselibrary{theorems,breakable}
\usepackage{comment}
\usepackage{needspace}
\usetikzlibrary{calc}
\usetikzlibrary{arrows.meta}

% רמז לסעיף: \hint{...} בשורה נפרדת, לפני ה-solbox של הסעיף.
% במאגר (\soltrue) הרמז מודפס בכחול; בחוברת לסטודנטים (\solfalse) הוא מוסתר, ובאתר הוא נפתח בלחיצה.
\ifsol
\newcommand{\hint}[1]{\par{\small\color{blue!55!black}\textbf{רמז:}\enskip #1}\par\smallskip}
\else
\newcommand{\hint}[1]{}
\fi

\ifsol
\newtcolorbox{solbox}{
    enhanced,
    breakable,
    colback=gray!8,
    borderline east={2pt}{0pt}{gray!80},
    borderline west={0.3pt}{0pt}{gray!80},
    borderline north={0.3pt}{0pt}{gray!80},
    borderline south={0.3pt}{0pt}{gray!80},
    frame hidden,
    sharp corners,
    fontupper=\small,
    before upper={\textbf{פיתרון:}\enskip},
    grow to left by=0cm,
    grow to right by=0cm}
\else
\excludecomment{solbox}
\fi

\newcommand{\lvec}{\overrightarrow}
\renewcommand{\labelitemi}{$\diamond$}

% שם הקובץ של השאלה, בקצה השמאלי של השורה שמעליה (משמש ב-bank.tex).
% את השם הזה כותבים ב-\input כדי לשלב את השאלה בדף תרגילים.
\newcommand{\qtag}[1]{{\footnotesize\color{gray}\hfill קובץ: \textenglish{\ttfamily #1}}\par\nopagebreak[4]\smallskip}
\newcommand{\bankq}[2]{\item \qtag{#1}\input{#2}}

% \R{...} (טקסט עברי בתוך נוסחה) מגיע מ-babel ואינו מוגדר ב-MotkePreamble
\providecommand{\R}[1]{\RL{#1}}

% הפניה לשאלה אחרת במאגר לפי המזהה שלה, למשל \qref{int-013}
\newcommand{\qref}[1]{\textenglish{#1}}

% שאלה שהוצעה כדי לסגור פער בין מאגר המבחנים לדפי התרגילים (נכתבה ע"י Claude, טרם אושרה).
% כל הקובץ עטוף בסביבה הזו, ולכן השאלה מופיעה בתיבה ורודה בכל מקום שמשלבים אותה.
% לאחר אישור -- מוחקים את \begin{suggested} ואת \end{suggested}, ואת השורה "% suggested: yes" מכותרת הקובץ.
% עד אז השאלה מופיעה במאגר (bank.pdf) בלבד, ולא בחוברת לסטודנטים ולא באתר.
\definecolor{suggestpink}{RGB}{214,51,132}
\newtcolorbox{suggested}{
    enhanced,
    breakable,
    colback=suggestpink!7,
    colframe=suggestpink!70,
    boxrule=0.6pt,
    arc=2pt,
    left=6pt, right=6pt, top=4pt, bottom=4pt,
    before skip=2pt,
    before upper={{\color{suggestpink}\small\textbf{הצעה חדשה -- שאלה שנכתבה כדי לסגור פער מול מאגר המבחנים (טרם אושרה)}}\par\smallskip}}

% סימון לפתרון שעדיין לא נבדק ואושר (נכתב ע"י Claude). לאחר בדיקה -- מוחקים את השורה
% \unverified מתוך ה-solbox ואת השורה "% verified: no" מכותרת הקובץ.
\newcommand{\unverified}{{\color{red}\textbf{[טיוטה -- פתרון שטרם נבדק]}}\par}

% סעיפים בשתי עמודות -- רק בגרסה ללא פתרונות (בחוברת לסטודנטים).
% בגרסה עם פתרונות הסעיפים נשארים בעמודה אחת, כדי שהפתרונות לא יידחסו.
% שימוש: \begin{enumerate}\bMs{2} \item ... \eMs\end{enumerate}
\setlength{\multicolsep}{4pt plus 2pt minus 1pt}
\newcommand{\bMs}[1]{\ifsol\else\begin{multicols}{#1}\fi}
\newcommand{\eMs}{\ifsol\else\end{multicols}\fi}

%%%%%%%%%%%%%%%%%%%%%%%%%%%%%%%%%%%%%%%%%%%%%%%%%%%%%%%%%%%%%%%%%%%%%%
% פקודות של הקורס בהסתברות. אותן פקודות מוגדרות לאתר ב-site/macros.json.
%%%%%%%%%%%%%%%%%%%%%%%%%%%%%%%%%%%%%%%%%%%%%%%%%%%%%%%%%%%%%%%%%%%%%%
\providecommand{\N}{} \renewcommand{\N}{\mathbb{N}}
\providecommand{\set}{} \renewcommand{\set}[1]{\left\{#1\right\}}
\providecommand{\setof}{} \renewcommand{\setof}[2]{\left\{\,#1 \;\middle|\; #2\,\right\}}
\providecommand{\symdiff}{} \renewcommand{\symdiff}{\mathbin{\triangle}}
\providecommand{\comp}{} \renewcommand{\comp}[1]{\overline{#1}}
\providecommand{\Pow}{} \renewcommand{\Pow}{\mathcal{P}}
\renewcommand{\emptyset}{\varnothing}
\providecommand{\PP}{} \renewcommand{\PP}{\mathbb{P}}
\providecommand{\E}{} \renewcommand{\E}{\mathbb{E}}
\providecommand{\Var}{} \renewcommand{\Var}{\operatorname{Var}}
\providecommand{\Ber}{} \renewcommand{\Ber}{\mathrm{Ber}}
\providecommand{\Bin}{} \renewcommand{\Bin}{\mathrm{Bin}}
\providecommand{\Geo}{} \renewcommand{\Geo}{\mathrm{Geo}}
\providecommand{\Poi}{} \renewcommand{\Poi}{\mathrm{Poi}}
% התשובה הסופית של סעיף, במסגרת
\providecommand{\ans}{} \renewcommand{\ans}[1]{\boxed{#1}}

% דיאגרמת ון לשלוש קבוצות A, B, C. הארגומנט הראשון מונה את האזורים הצבועים: 1 = בתוך הקבוצה,
% 0 = מחוצה לה (לפי הסדר A/B/C). למשל \venn{1/1/0,0/1/1}{(ג)}. באתר הדיאגרמה מצוירת ב-SVG.
\definecolor{venncolor}{RGB}{35,72,168}
\newcommand{\vclip}[2]{\ifnum#2=1\relax\begin{scope}\clip #1;\else
  \begin{scope}[even odd rule]\clip (-1.6,-2.75) rectangle (3.0,1.65) #1;\fi}
\newcommand{\vennatom}[3]{\vclip{\vA}{#1}\vclip{\vB}{#2}\vclip{\vC}{#3}%
  \fill[venncolor!30] (-1.6,-2.75) rectangle (3.0,1.65);\end{scope}\end{scope}\end{scope}}
\newcommand{\venn}[2]{%
\begin{tikzpicture}[scale=0.42,baseline=(current bounding box.center)]
  \def\vA{(0,0) circle (1.1)}\def\vB{(1.4,0) circle (1.1)}\def\vC{(0.7,-1.15) circle (1.1)}
  \foreach \a/\b/\c in {#1} {\vennatom{\a}{\b}{\c}}
  \draw (-1.6,-2.75) rectangle (3.0,1.65);
  \draw[thick]\vA; \draw[thick]\vB; \draw[thick]\vC;
  \node at (-0.55,0.55) {\scriptsize $A$}; \node at (1.95,0.55) {\scriptsize $B$}; \node at (0.7,-1.85) {\scriptsize $C$};
  \node[below] at (0.7,-2.75) {\small #2};
\end{tikzpicture}}

\newcommand{\bookletkind}{רמזים ופתרונות}
\newcommand{\bookletsitenote}{באתר התרגול הרמזים והפתרונות נפתחים בלחיצה, שלב אחר שלב.}
%%%%%%%%%%%%%%%%%%%%%%%%%%%%%%%%%%%%%%%%%%%%%%%%%%%%%%%%%%%%%%%%%%%%%%
% עמוד השער של חוברת התרגילים לסטודנטים (booklet.tex).
% ערכו כאן את הסמסטר, שם הקורס וההערה -- booklet.tex עצמו נוצר אוטומטית.
%%%%%%%%%%%%%%%%%%%%%%%%%%%%%%%%%%%%%%%%%%%%%%%%%%%%%%%%%%%%%%%%%%%%%%
\newcommand{\semester}{סמסטר א' תשפ״ז}
\newcommand{\course}{הסתברות לסטטיסטיקה}
\newcommand{\bookletnote}{%
החוברת מרכזת תרגילים על כל נושאי הקורס, ותעודכן במהלך הסמסטר.
}

% כתובת אתר התרגול (רמזים ופתרונות שלב אחר שלב). מופיעה בעמוד השער, עם קוד QR.
% \string_ : a literal underscore, both in the printed address and in the QR code
\edef\siteurl{https://motkep173.github.io/Ex2027A\string_9921/}
% הכותרת הגדולה בעמוד השער. חוברת הרמזים והפתרונות (solutions.tex) מחליפה אותה.
\providecommand{\bookletkind}{חוברת תרגילים}
\providecommand{\bookletsitenote}{רמזים ופתרונות מלאים, שנפתחים שלב אחר שלב, נמצאים באתר התרגול.
אפשר גם להוריד משם את חוברת הרמזים והפתרונות.}
\usepackage{qrcode}

\newcommand{\booklettitlepage}{%
\begin{titlepage}
\begin{textblock*}{9cm}[0.5,0.5](105mm,4cm)
	\centering\includegraphics[width=2.2in]{preamble/logo.jpg}
\end{textblock*}
\vspace*{5cm}
\begin{center}
{\large היחידה להוראת קורסי יסוד}\\[0.8em]
{\Large\bf \course\ -- 214.1.9921}\\[3em]
\textcolor{bkorange}{\fontsize{34}{40}\selectfont\bfseries \bookletkind}\\[1em]
\textcolor{bkgray!60}{\rule{0.5\textwidth}{1.2pt}}\\[1.5em]
\textcolor{bkgray}{\Large \semester}\\[4em]
\begin{minipage}{0.75\textwidth}
\centering\large \bookletnote
\end{minipage}
\vfill
\begin{minipage}{0.62\textwidth}
\raggedleft\bookletsitenote\\[0.4em]
{\small\textenglish{\ttfamily \siteurl}}
\end{minipage}\hspace{1.2em}%
\begin{minipage}{2.6cm}
\centering\qrcode[height=2.4cm]{\siteurl}
\end{minipage}\par
\vspace{2.5em}
{\small עודכן לאחרונה: \today}
\end{center}
\end{titlepage}
}

%%%%%%%%%%%%%%%%%%%%%%%%%%%%%%%%%%%%%%%%%%%%%%%%%%%%%%%%%%%%%%%%%%%%%%
% עיצוב החוברת: צבעים, כותרות הפרקים והפרדה בין התרגילים.
%%%%%%%%%%%%%%%%%%%%%%%%%%%%%%%%%%%%%%%%%%%%%%%%%%%%%%%%%%%%%%%%%%%%%%
\definecolor{bkorange}{HTML}{E46C0A}
\definecolor{bkgray}{HTML}{5A5A5A}
\definecolor{bklight}{HTML}{F3F3F3}

% סגנון ההפרדה בין התרגילים: lines (קו אפור), boxes (תיבה אפורה), space (מרווח בלבד)
\providecommand{\bookletstyle}{lines}

% כותרת פרק: מספר ושם בכתום, וקו אפור מתחת
\makeatletter
\renewcommand{\section}{\@ifstar\bk@sectionstar\bk@section}
\newcommand{\bk@heading}[1]{%
  {\noindent\color{bkorange}\Huge\bfseries #1\par}\nobreak\vspace{4pt}%
  {\color{bkgray!50}\hrule height 1.2pt}\nobreak\vspace{1.2em}}
\newcommand{\bk@sectionstar}[1]{\bk@heading{#1}\par\@afterheading}
\newcommand{\bk@section}[1]{%
  \refstepcounter{section}%
  \addcontentsline{toc}{section}{\protect\numberline{\thesection}#1}%
  \bk@heading{\thesection\quad #1}}
\makeatother

% תרגיל אחד בחוברת: רשימה עם פריט אחד, שהמספר שלו הוא המונה bkq (רץ מתחילת החוברת).
% כך הסעיפים (א), (ב)... נשארים ברמה השנייה של הרשימות, כמו קודם.
% \bookq[כותרת פרק]{קובץ}: התרגיל לא נשבר בין עמודים. לפני הדפסתו מודדים את גובהו (כולל כותרת הפרק
% שלפניו, אם יש), ואם אין לו מספיק מקום בעמוד הנוכחי הוא עובר לעמוד הבא.
% תרגיל עם פתרון שארוך מעמוד שלם: נשמר יחד לפחות נוסח השאלה (בלי הפתרון) ותחילת הפתרון.
% פרק חדש לא מתחיל בעמוד חדש: הכותרת שלו נשארת עם תחילת התרגיל הראשון בפרק, והתרגיל הזה יכול להמשיך בעמוד הבא.
\makeatletter
\edef\bk@style{\bookletstyle}
\def\bk@lines{lines}
\def\bk@boxes{boxes}
\newcounter{bkq}
\newcommand{\bknum}[1]{\textcolor{bkorange}{\bfseries #1.}}
\newcommand{\bk@q}[1]{\stepcounter{bkq}%
  \begin{enumerate}[label=\bknum{\arabic{bkq}}, leftmargin=2em, topsep=0pt]\item \input{#1}\end{enumerate}}
\newtcolorbox{bkbox}{enhanced, breakable, colback=bklight, frame hidden, sharp corners,
  borderline east={2.5pt}{0pt}{bkorange}, left=4pt, right=6pt, top=5pt, bottom=5pt,
  before skip=0pt, after skip=10pt}
\ifx\bk@style\bk@boxes
  \newcommand{\bk@one}[1]{\begin{bkbox}\bk@q{#1}\end{bkbox}}
\else\ifx\bk@style\bk@lines
  \newcommand{\bk@one}[1]{\bk@q{#1}\par\vspace{8pt}{\color{bkgray!35}\hrule height 0.5pt}\vspace{8pt}}
\else
  \newcommand{\bk@one}[1]{\bk@q{#1}\par\vspace{16pt}}
\fi\fi
% מדידה: מסדרים את התרגיל בתיבה שלא מודפסת, ומחזירים את המונים למצבם.
% (הכתיבה לקבצי aux מתוך התיבה לא מתבצעת, כי התיבה לא נשלחת לעמוד.)
\newbox\bk@measure
\newcount\bk@cq \newcount\bk@ceq \newcount\bk@csec
\newdimen\bk@need
\newcommand{\bk@measureit}[1]{%
  \bk@cq=\value{bkq}\bk@ceq=\value{equation}\bk@csec=\value{section}%
  \setbox\bk@measure=\vbox{\hsize\linewidth #1}%
  \setcounter{bkq}{\bk@cq}\setcounter{equation}{\bk@ceq}\setcounter{section}{\bk@csec}%
  \bk@need=\dimexpr\ht\bk@measure+\dp\bk@measure\relax}
\newcommand{\bookq}[2][]{%
  \par
  \if\relax\detokenize{#1}\relax
    \bk@keep{#2}%
  \else
    % תחילת פרק: הפרק ממשיך באותו עמוד אם יש מקום לכותרת ולתחילת התרגיל הראשון,
    % גם אם התרגיל הזה ימשיך בעמוד הבא
    \needspace{7\baselineskip}%
    #1\bk@one{#2}%
  \fi}
\newcommand{\bk@keep}[1]{%
  \bk@measureit{\bk@one{#1}}%
  \ifdim\bk@need>\dimexpr\textheight-3\baselineskip\relax
    % ארוך מעמוד: מודדים רק את נוסח השאלה, בלי הפתרונות
    \bk@measureit{\RenewEnviron{solbox}{}\bk@one{#1}}%
    \advance\bk@need by 4\baselineskip
    \ifdim\bk@need>\textheight \bk@need=\textheight \fi
  \fi
  \needspace{\bk@need}%
  \bk@one{#1}}
\makeatother

% פתרון שעדיין לא נבדק מסומן כך, כמו באתר
\renewcommand{\unverified}{{\color{bkgray}\small\textbf{פתרון שלא נבדק על ידי צוות הקורס}}\par}

\begin{document}
\selectlanguage{hebrew}
\booklettitlepage
{\small\tableofcontents}

\expandafter\def\csname qref@sets-005\endcsname{1}
\expandafter\def\csname qref@sets-008\endcsname{2}
\expandafter\def\csname qref@sets-001\endcsname{3}
\expandafter\def\csname qref@sets-002\endcsname{4}
\expandafter\def\csname qref@sets-003\endcsname{5}
\expandafter\def\csname qref@sets-006\endcsname{6}
\expandafter\def\csname qref@sets-007\endcsname{7}
\expandafter\def\csname qref@sets-011\endcsname{8}
\expandafter\def\csname qref@sets-010\endcsname{9}
\expandafter\def\csname qref@sets-004\endcsname{10}
\expandafter\def\csname qref@sets-009\endcsname{11}
\expandafter\def\csname qref@comb-001\endcsname{12}
\expandafter\def\csname qref@comb-015\endcsname{13}
\expandafter\def\csname qref@comb-016\endcsname{14}
\expandafter\def\csname qref@comb-019\endcsname{15}
\expandafter\def\csname qref@comb-021\endcsname{16}
\expandafter\def\csname qref@comb-002\endcsname{17}
\expandafter\def\csname qref@comb-008\endcsname{18}
\expandafter\def\csname qref@comb-009\endcsname{19}
\expandafter\def\csname qref@comb-010\endcsname{20}
\expandafter\def\csname qref@comb-014\endcsname{21}
\expandafter\def\csname qref@comb-018\endcsname{22}
\expandafter\def\csname qref@comb-005\endcsname{23}
\expandafter\def\csname qref@comb-006\endcsname{24}
\expandafter\def\csname qref@comb-003\endcsname{25}
\expandafter\def\csname qref@comb-011\endcsname{26}
\expandafter\def\csname qref@comb-012\endcsname{27}
\expandafter\def\csname qref@comb-013\endcsname{28}
\expandafter\def\csname qref@comb-022\endcsname{29}
\expandafter\def\csname qref@comb-027\endcsname{30}
\expandafter\def\csname qref@comb-028\endcsname{31}
\expandafter\def\csname qref@comb-023\endcsname{32}
\expandafter\def\csname qref@comb-024\endcsname{33}
\expandafter\def\csname qref@comb-029\endcsname{34}
\expandafter\def\csname qref@prob-006\endcsname{35}
\expandafter\def\csname qref@prob-007\endcsname{36}
\expandafter\def\csname qref@prob-015\endcsname{37}
\expandafter\def\csname qref@prob-016\endcsname{38}
\expandafter\def\csname qref@prob-017\endcsname{39}
\expandafter\def\csname qref@prob-002\endcsname{40}
\expandafter\def\csname qref@prob-010\endcsname{41}
\expandafter\def\csname qref@prob-012\endcsname{42}
\expandafter\def\csname qref@comb-004\endcsname{43}
\expandafter\def\csname qref@prob-001\endcsname{44}
\expandafter\def\csname qref@prob-004\endcsname{45}
\expandafter\def\csname qref@prob-005\endcsname{46}
\expandafter\def\csname qref@prob-008\endcsname{47}
\expandafter\def\csname qref@prob-009\endcsname{48}
\expandafter\def\csname qref@prob-011\endcsname{49}
\expandafter\def\csname qref@prob-013\endcsname{50}
\expandafter\def\csname qref@prob-014\endcsname{51}
\renewcommand{\qref}[1]{\ifcsname qref@#1\endcsname\csname qref@#1\endcsname\else\textenglish{#1}\fi}

\newpage
\section*{רמזים}
\addcontentsline{toc}{section}{רמזים}
כאן מופיעים רק הרמזים, בלי הפתרונות. נסו קודם לפתור בעזרת הרמז, ורק אחר כך עברו לפתרון.


\subsection*{תורת הקבוצות ומאורעות}
\noindent \bknum{7}\ (א)\enskip לכל אזור צבוע, בדקו אם הוא בתוך כל אחת מהקבוצות $A,B,C$ או מחוץ לה.\par\smallskip
\noindent \bknum{8}\ (א)\enskip הכלה דו-כיוונית: הראו ש-$A\cap B\subseteq A$ וגם ש-$A\subseteq A\cap B$. בכיוון השני השתמשו בהנחה $A\subseteq B$.\par\smallskip
\noindent \bknum{9}\enskip כתבו את $A\cup B$ כאיחוד של שתי קבוצות זרות, $A\setminus B$ ו-$B$. לגודל של $A\setminus B$ השתמשו בעקרון החיסור.\par\smallskip
\noindent \bknum{11}\ (ב)\enskip תתי-הקבוצות של קבוצה עם איבר אחד הן הקבוצה הריקה והקבוצה עצמה.\par\smallskip

\subsection*{קומבינטוריקה}
\noindent \bknum{12}\ (ג)\enskip חשבו על כל נושא כעל "בלוק" אחד: קודם מסדרים את הבלוקים, ואחר כך את הספרים בתוך כל בלוק.\par\smallskip
\noindent \bknum{12}\ (ה)\enskip בכל שבוע בוחרים זוג ספרים, בלי חשיבות לסדר בתוך הזוג. הרשימה היא סדרה של 5 זוגות כאלה, אחד לכל שבוע, ולכן אפשר להשתמש בעקרון הכפל.\par\smallskip
\noindent \bknum{14}\ (ה)\enskip במספר כזה יש ספרה $X$ שמופיעה 3 פעמים, וספרה אחרת $Y$ שמופיעה פעם אחת. כמה אפשרויות יש לבחירת $X$, לבחירת $Y$, ולמקום של $Y$?\par\smallskip
\noindent \bknum{15}\ (ב)\enskip כמה סידורים של ספרים שונים נראים בדיוק אותו דבר, כשמסתכלים רק על הצבעים?\par\smallskip
\noindent \bknum{16}\ (ג)\enskip חשבו על שתי חנויות ההלבשה וחנות הנעליים שביניהן כעל "בלוק" של 3 מקומות רצופים.\par\smallskip
\noindent \bknum{22}\ (ד)\enskip קל יותר לספור את המשלים: בכמה דרכים קומה 4 \emph{לא} מיוצגת?\par\smallskip
\noindent \bknum{23}\ (א)\enskip כתבו את 6 כשורה של שש אחדות: $1\ \_\ 1\ \_\ 1\ \_\ 1\ \_\ 1\ \_\ 1$. בין האחדות יש 5 רווחים, ובכל רווח מחליטים אם לשים בו "$+$" או להשאיר אותו ריק.\par\smallskip
\noindent \bknum{25}\enskip גפן צריכה לארוז 13 זוגות מכנסיים לחופשה. בארון יש לה 30 זוגות, מתוכם 10 ארוכים ו-20 קצרים. בכמה דרכים היא יכולה לעשות זאת?\par\smallskip
\noindent \bknum{26}\enskip סמנו $|\Omega|=n$, ובחרו ערכים מתאימים ל-$a$ ול-$b$ בבינום של ניוטון.\par\smallskip
\noindent \bknum{27}\enskip כמה זוגות מכנסיים יש עכשיו לגפן בארון, וכמה היא צריכה לארוז?\par\smallskip
\noindent \bknum{28}\ (א)\enskip מתוך $n$ תלמידים, שאחד מהם הוא דני, בוחרים $k$. מפרידים לשני מקרים: דני נבחר, או דני לא נבחר.\par\smallskip
\noindent \bknum{29}\ (ב)\enskip אם ב-$n$ המקומות הראשונים יש $i$ אפסים, כמה אפסים יש בסדרה כולה?\par\smallskip
\noindent \bknum{30}\ (ד)\enskip קל יותר לספור את הקבוצות שבהן שניהם נמצאים.\par\smallskip
\noindent \bknum{31}\enskip כמה תתי-קבוצות \emph{לא} מקיימות את התנאי?\par\smallskip
\noindent \bknum{32}\ (ז)\enskip כל מחובר בפיתוח הוא מהצורה $\binom6k(4x)^k3^{6-k}$.\par\smallskip
\noindent \bknum{32}\ (ח)\enskip כל מחובר בפיתוח הוא מהצורה $\binom5k(a^2)^k(b^3)^{5-k}$. לאיזה $k$ מקבלים $a^6b^6$?\par\smallskip
\noindent \bknum{33}\ (א)\enskip מתוך $n$ סטודנטים בוחרים ועדה של $r$ חברים, ומתוך הוועדה בוחרים ועדת משנה של $k$ חברים. בכמה דרכים אפשר לעשות את זה, אם מתחילים מהוועדה? ואם מתחילים מוועדת המשנה?\par\smallskip
\noindent \bknum{34}\ (א)\enskip מתוך $n$ סטודנטים בוחרים ועדה (בכל גודל), ובוחרים לה יו"ר מבין חבריה.\par\smallskip
\noindent \bknum{34}\ (ב)\enskip בכיתה $n$ בנים ו-$n$ בנות, ובוחרים שני תלמידים.\par\smallskip

\subsection*{מבוא להסתברות}
\noindent \bknum{38}\enskip קודם בחרו את שני החודשים. בכמה דרכים אפשר לחלק את 6 האנשים בין שני החודשים, כך שאף חודש לא יישאר ריק?\par\smallskip
\noindent \bknum{39}\ (ב)\enskip חשבו על שני הספרים כעל „גוש” אחד.\par\smallskip
\noindent \bknum{40}\ (ב)\enskip האם המאורעות $\set{1,2},\set{2,3},\set{3,4},\dots$ זרים? חפשו מאורעות מהצורה $\set{k,k+1}$ שכן זרים בזוגות, ושהאיחוד שלהם הוא $\Omega$.\par\smallskip
\noindent \bknum{44}\ (ג)\enskip סמנו $A_i$ = "לסניף ה-$i$ הגיעה לפחות מדענית מחשב אחת", ועברו למשלים: $\PP(A_1\cap A_2\cap A_3)=1-\PP(A_1^C\cup A_2^C\cup A_3^C)$.\par\smallskip
\noindent \bknum{47}\ (א)\enskip סמנו $A$ = „עבר את הראשון”, $B$ = „עבר את השני”. אפשר גם לסדר את הנתונים בטבלה של $2\times2$: עבר/נכשל בראשון, ועבר/נכשל בשני.\par\smallskip

\newpage
\section*{פתרונות}
\addcontentsline{toc}{section}{פתרונות}
\renewcommand{\hint}[1]{}   % the hints were listed above
\setcounter{section}{0}


\bookq[\section{תורת הקבוצות ומאורעות}]{questions/01-sets/sets-005}
\bookq[]{questions/01-sets/sets-008}
\bookq[]{questions/01-sets/sets-001}
\bookq[]{questions/01-sets/sets-002}
\bookq[]{questions/01-sets/sets-003}
\bookq[]{questions/01-sets/sets-006}
\bookq[]{questions/01-sets/sets-007}
\bookq[]{questions/01-sets/sets-011}
\bookq[]{questions/01-sets/sets-010}
\bookq[]{questions/01-sets/sets-004}
\bookq[]{questions/01-sets/sets-009}


\bookq[\section{קומבינטוריקה}]{questions/02-combinatorics/comb-001}
\bookq[]{questions/02-combinatorics/comb-015}
\bookq[]{questions/02-combinatorics/comb-016}
\bookq[]{questions/02-combinatorics/comb-019}
\bookq[]{questions/02-combinatorics/comb-021}
\bookq[]{questions/02-combinatorics/comb-002}
\bookq[]{questions/02-combinatorics/comb-008}
\bookq[]{questions/02-combinatorics/comb-009}
\bookq[]{questions/02-combinatorics/comb-010}
\bookq[]{questions/02-combinatorics/comb-014}
\bookq[]{questions/02-combinatorics/comb-018}
\bookq[]{questions/02-combinatorics/comb-005}
\bookq[]{questions/02-combinatorics/comb-006}
\bookq[]{questions/02-combinatorics/comb-003}
\bookq[]{questions/02-combinatorics/comb-011}
\bookq[]{questions/02-combinatorics/comb-012}
\bookq[]{questions/02-combinatorics/comb-013}
\bookq[]{questions/02-combinatorics/comb-022}
\bookq[]{questions/02-combinatorics/comb-027}
\bookq[]{questions/02-combinatorics/comb-028}
\bookq[]{questions/02-combinatorics/comb-023}
\bookq[]{questions/02-combinatorics/comb-024}
\bookq[]{questions/02-combinatorics/comb-029}


\bookq[\section{מבוא להסתברות}]{questions/03-probability/prob-006}
\bookq[]{questions/03-probability/prob-007}
\bookq[]{questions/03-probability/prob-015}
\bookq[]{questions/03-probability/prob-016}
\bookq[]{questions/03-probability/prob-017}
\bookq[]{questions/03-probability/prob-002}
\bookq[]{questions/03-probability/prob-010}
\bookq[]{questions/03-probability/prob-012}
\bookq[]{questions/02-combinatorics/comb-004}
\bookq[]{questions/03-probability/prob-001}
\bookq[]{questions/03-probability/prob-004}
\bookq[]{questions/03-probability/prob-005}
\bookq[]{questions/03-probability/prob-008}
\bookq[]{questions/03-probability/prob-009}
\bookq[]{questions/03-probability/prob-011}
\bookq[]{questions/03-probability/prob-013}
\bookq[]{questions/03-probability/prob-014}

\end{document}
